% Programming with InitechBase, Chapter 13: Testing and Debugging (AN-16;
% annex Sec 8.2 P15).  Panels replayed (examples/p13-*.dot); the program
% BADTOTAL.PRG; the memory variables of a main program remain after it
% (cmd/flow.c level 0, observed).
\chapter{Testing and Debugging}
\Gates{test-pubs-examples, test-interp-proc}
A program that runs is not yet a program that is right.  This chapter gives the
habits that find mistakes early, and shows how to find the command at fault
when a program stops.

\begin{ChapterCovers}
\item Testing on copies
\item Looking at a program after it ends
\item Finding the command at fault
\item Testing a part of a program at the dot prompt
\end{ChapterCovers}

\ATEHead{Preparing for This Chapter}
Start InitechBase with a fresh copy of the sample diskette.

\ATEHead{Testing on Copies}
Test a program on a copy of the databases, never on the files of the business.
A program that changes records -- \file{NEWINV.PRG}, for one -- writes to the disk
at once.  Keep a copy of the files as they were before the test, so that you can
run it again from the same start.

\ATEHead{Looking After the Run}
The memory variables a program creates remain after it ends, and so does the
record pointer.  Look at them with \cmd{?}:
\DotExample{examples/p13-check}
\fn{TYPE()} tells you whether a memory variable has the type you meant; a total
that turns out to be of type C was given a character value somewhere.

\ATEHead{Finding the Command at Fault}
\file{BADTOTAL.PRG}, listed in Chapter~12, stops with a message:
\DotExample{examples/p13-bad}
The message says that a name was not found.  \texttt{MTOTAL} is still 0 and the
record pointer is on record 1: the program stopped inside the loop, on the
first record, before anything was added.  The only name used there that is not
a memory variable is \texttt{AMOUNTS}: the field is \texttt{AMOUNT}.  Correct the
line with your editor and run the program again.

\ATEHead{Testing a Part at the Dot Prompt}
Every command of a program can be typed at the dot prompt, except the
structures \cmd{IF}, \cmd{DO CASE} and \cmd{DO WHILE}.  To test a calculation,
open the database, move to a record, store the memory variables the
calculation needs, and type the calculation with \cmd{?}.  When it gives the
right result, copy it into the program.

\begin{Tip}
Add a \cmd{?} to a program to display a memory variable at the point where you
suspect the mistake, and remove it when the program is right.
\end{Tip}

\ATEHead{Where Do You Go From Here?}
Chapter~14 works with two databases at once.

\ATEHead{Summary}
\begin{enumerate}
\item Test on copies of the databases.
\item Memory variables and the record pointer remain after a program; look at
  them with \cmd{?} and \fn{TYPE()}.
\item The message, the record number and the memory variables lead you to the
  command at fault.
\item Try calculations at the dot prompt before you put them in a program.
\end{enumerate}
